% 20260707_Assingment_Algebra_IIA
\documentclass[dvipdfmx]{jsreport}
\usepackage{soupreamble}

\title{代数学特論IIA 第3回レポート課題}
\author{M1 6012 川村聡一郎}
\date{}


\begin{document}
\maketitle

\begin{tcolorbox}
	$F(z)$が開領域$D$で(定数でない)有理関数で$\partial D$で零点も極も持たないものとする．$D$内の$F(z)$の零点の個数を$N$個，極の個数を$M$個とおくと
\begin{equation*}
	\frac{1}{2\pi i}\int_{\partial\mathcal{D}}\frac{F'(z)}{F(z)}dz= N- M
\end{equation*}
	であることを示せ．
\end{tcolorbox}

\begin{souanswer} %答案はここに書く
	$a_j$を$F(z)$の零点とすると，
\begin{equation}
	F(z)= \prod_{j= 1}^N (z- a_j)\cdot h(z),\ \ h(z)\ne 0
\end{equation}
	と書ける．これを$z$で微分して$F(z)$で割ると，
\begin{equation}
	\frac{F'(z)}{F(z)}= \sum_{j= 1}^N \frac{1}{z- a_j}+ \frac{h'(z)}{h(z)}
\end{equation}
	とかける．ここで$F'(z)/ F(z)$を$\partial D$まわりで周回積分すると，
\begin{equation}
	\int_{\partial D}\frac{F'(z)}{F(z)}dz
	= \int_{\partial D}\left(\sum_{j= 1}^N \frac{1}{z- a_j}+ \frac{h'(z)}{h(z)}\right)dz
	= 2\pi i\sum_{j= 1}^N \Res{\left(\frac{1}{z- a_j}, a_j\right)}+ 0
	= 2\pi iN
\end{equation}
	次に極について考える．$F(z)$の極を$b_k\ \ (k= 1, 2, \ldots, M)$とおくと，
\begin{equation}
	F(z)= \frac{1}{(z- b_1)(z- b_2)\cdots (z- b_M)}\cdot g(z)= \left(\prod_{k= 1}^M (z- b_k)^{-1}\right)\cdot g(z),\ \ g(z)\ne 0
\end{equation}
	同様に$z$で微分して$F(z)$で割ると，
\begin{equation}
	\frac{F'(z)}{F(z)}= -\sum_{k= 1}^M (z- b_k)^{-1}+ \frac{g'(z)}{g(z)}
\end{equation}
ここで$F'(z)/ F(z)$を$\partial D$まわりで周回積分すると，
\begin{equation}
	\int_{\partial D}\frac{F'(z)}{F(z)}dz
	= \int_{\partial D}\left(-\sum_{k= 1}^M (z- b_k)^{-1}+ \frac{g'(z)}{g(z)}\right)dz
	= 2\pi i\sum_{j= 1}^M \Res{\left(\frac{1}{z- b_k}, b_k\right)}+ 0
	= -2\pi iM
\end{equation}
	$F(z)$は零点も極も持つ関数だから
\begin{equation}
	F(z)= \frac{\displaystyle \prod_{j= 1}^N (z- a_j)}{\displaystyle \prod_{k= 1}^M (z- b_k)}\cdot p(z),\ \ p(z)\ne 0
\end{equation}
	とかけ，$F'(z)/F(z)$の積分の留数は
\begin{equation}
	\frac{1}{2\pi i}\int_{\partial D} \frac{F'(z)}{F(z)}dz
	= N- M
\end{equation}
\end{souanswer}
\end{document}